Put \(q=m-2c=h-c\). By \(\text{def:affine-separation}\), if \(q<3\), then \(\gamma_H=0\), contrary to \(\gamma_H\geq\gamma_0>0\). Hence
\[
q\geq3.
\tag{1}
\]

\textbf{Uniform structural bounds.}
Let \(G\) be an invertible matrix with unit diagonal, \(\operatorname{CP}(I_p-G)\leq1\), and
\(\|G\|_{\mathrm{op}}\|G^{-1}\|_{\mathrm{op}}\leq\overline\kappa\). In the Leibniz expansion of \(\det G\), a permutation is a disjoint union of fixed points and directed simple cycles. Unit diagonal makes every fixed-point factor one, while every nontrivial-cycle factor has absolute value at most \(\operatorname{CP}(I_p-G)\leq1\). Thus every Leibniz monomial has absolute value at most one and
\[
|\det G|\leq p!.
\tag{2}
\]
If \(\sigma_1(G)\geq\cdots\geq\sigma_p(G)>0\) are the singular values, the condition-number bound gives \(\sigma_j(G)\geq\sigma_1(G)/\overline\kappa\) for \(j\geq2\). Since the product of the singular values is \(|\det G|\), (2) yields
\[
\frac{\sigma_1(G)^p}{\overline\kappa^{p-1}}
\leq\prod_{j=1}^p\sigma_j(G)
=|\det G|
\leq p!,
\]
and therefore
\[
\|G\|_{\mathrm{op}}
\leq\bigl(p!\,\overline\kappa^{p-1}\bigr)^{1/p}
=L_0.
\tag{3}
\]
Because a unit diagonal entry has absolute value one, \(\|G\|_{\mathrm{op}}\geq1\); hence
\[
\|G^{-1}\|_{\mathrm{op}}\leq\overline\kappa.
\tag{4}
\]
The true matrix satisfies the hypotheses because \(\zeta_D\geq\zeta_0>0\), and every feasible candidate and every matrix in its one-sided normalization closure satisfies them with \(\operatorname{CP}(I_p-G)\leq1\). Thus (3)--(4) apply on both sides without a uniform candidate normalization slack.

\textbf{Compact comparison family.}
Define
\[
\rho_0=\min\left\{1,\frac{1}{2K_0\overline M}\right\}>0.
\tag{5}
\]
For each \(H_1\subseteq E\) with \(|H_1|=h\) and each \(S\subseteq H_1\) with \(|S|=q\), let the corresponding comparison family consist of tuples
\[
\tau=\bigl(D,U,\Omega,(s_e,\Sigma_e)_{e\in H_1},
A,V,\Psi,(t_e,\Gamma_e)_{e\in S}\bigr)
\tag{6}
\]
with
\[
DU=UD=I_p,
\qquad AV=VA=I_p,
\tag{7}
\]
unit diagonals for \(D,A\), and the one-sided normalization conditions
\[
\operatorname{CP}(I_p-D)\leq1-\zeta_0,
\qquad
\operatorname{CP}(I_p-A)\leq1,
\tag{8}
\]
together with
\[
\|D\|_{\mathrm{op}}\|U\|_{\mathrm{op}}\leq\overline\kappa,
\qquad
\|A\|_{\mathrm{op}}\|V\|_{\mathrm{op}}\leq\overline\kappa.
\tag{9}
\]
Require also
\[
\Omega,\Sigma_e\in\mathbb S_+^p,\quad s_e\in[0,\infty)^p
\quad(e\in H_1),
\]
\[
\Psi,\Gamma_e\in\mathbb S_+^p,\quad t_e\in[0,\infty)^p
\quad(e\in S),
\tag{10}
\]
and
\[
D\Sigma_eD^{\mathsf T}=\Omega+\operatorname{diag}(s_e)
\quad(e\in H_1),
\qquad
A\Gamma_eA^{\mathsf T}=\Psi+\operatorname{diag}(t_e)
\quad(e\in S).
\tag{11}
\]
Finally require that the affine-separation functional computed from \((H_1,s)\) is at least \(\gamma_0\), that
\[
1+\|\Omega\|_{\mathrm{op}}
+\max_{e\in H_1}\|\Sigma_e\|_{\mathrm{op}}
+\max_{e\in H_1}\|\operatorname{diag}(s_e)\|_{\mathrm{op}}
\leq\overline M,
\tag{12}
\]
and that
\[
\|\Gamma_e-\Sigma_e\|_{\mathrm{op}}\leq1
\quad(e\in S).
\tag{13}
\]
Let \(\mathfrak T\) be the finite disjoint union over \((H_1,S)\).

Equations (3)--(4) bound \(D,U,A,V\). Equation (12) gives \(\|\Sigma_e\|_{\mathrm{op}}\leq\overline M-1\), so (13) gives
\[
\|\Gamma_e\|_{\mathrm{op}}
\leq\|\Sigma_e\|_{\mathrm{op}}
+\|\Gamma_e-\Sigma_e\|_{\mathrm{op}}
\leq\overline M.
\tag{14}
\]
By (3), (11), and (14),
\[
\|\Psi+\operatorname{diag}(t_e)\|_{\mathrm{op}}
=\|A\Gamma_eA^{\mathsf T}\|_{\mathrm{op}}
\leq L_0^2\overline M.
\tag{15}
\]
Both summands on the left are positive semidefinite. For positive-semidefinite \(P,Q\), the Rayleigh-quotient inequalities \(0\leq v^{\mathsf T}Pv\leq v^{\mathsf T}(P+Q)v\) for every unit vector \(v\) imply \(\|P\|_{\mathrm{op}}\leq\|P+Q\|_{\mathrm{op}}\), and likewise for \(Q\). Consequently,
\[
\|\Psi\|_{\mathrm{op}}\leq L_0^2\overline M,
\qquad
\|\operatorname{diag}(t_e)\|_{\mathrm{op}}
\leq L_0^2\overline M.
\tag{16}
\]
Thus every coordinate in (6) is bounded. Each fixed-index family is closed: (7) and (11) are polynomial equalities; the cones and orthants in (10) are closed; the norm inequalities are closed; \(\operatorname{CP}\) is the maximum of finitely many continuous absolute cycle-product monomials; and the true affine-separation functional is a finite minimum of finite maxima of continuous absolute determinants. The weak candidate condition in (8) is precisely what closes the normalization side. Hence \(\text{lem:euclidean-heine-borel}\) makes each fixed-index family compact, and the finite disjoint union \(\mathfrak T\) is compact.

Define the continuous residual
\[
\mathfrak r(\tau)=
\max_{e\in S}
\left\|A\Sigma_eA^{\mathsf T}
-\Psi-\operatorname{diag}(t_e)\right\|_{\mathrm F}
\tag{17}
\]
and the closed subfamily
\[
\mathfrak T_{\mathrm{far}}
=\{\tau\in\mathfrak T:\|AU-I_p\|_{\mathrm F}\geq\rho_0\}.
\tag{18}
\]
It is compact. If it is empty, set \(\mu_0=1\). Otherwise, \(\text{lem:compact-extreme-value}\) gives an attaining minimum
\[
\mu_0=\min_{\tau\in\mathfrak T_{\mathrm{far}}}\mathfrak r(\tau).
\tag{19}
\]

We show \(\mu_0>0\). Suppose, to the contrary in the nonempty case, that an attaining tuple has \(\mathfrak r(\tau)=0\). Then for every \(e\in S\),
\[
D\Sigma_eD^{\mathsf T}=\Omega+\operatorname{diag}(s_e),
\qquad
A\Sigma_eA^{\mathsf T}=\Psi+\operatorname{diag}(t_e).
\tag{20}
\]
Set \(Q=AU=AD^{-1}\), fix \(e_0\in S\), and put \(x_e=s_e-s_{e_0}\). Subtracting (20) at \(e_0\) from (20) at \(e\) cancels both invariant nuisances and gives
\[
Q\operatorname{diag}(x_e)Q^{\mathsf T}
\text{ is diagonal for every }e\in S.
\tag{21}
\]
Let \(W=\operatorname{span}\{x_e:e\in S\}\). By linearity, (21) holds with \(x_e\) replaced by any \(w\in W\). For every \(k<\ell\), the coordinate projection of \(W\) onto coordinates \((k,\ell)\) equals \(\mathbb R^2\). Indeed, otherwise a nonzero \((u,v)\) would annihilate that projection, placing every point \((s_{ek},s_{e\ell})\), \(e\in S\), on one affine line through the point indexed by \(e_0\). This contradicts the affine-separation bound on the particular \(q\)-element subset \(S\), which supplies a triple with determinant of absolute value at least \(\gamma_0>0\).

Choose a basis \(w^1,\ldots,w^d\) of \(W\) and set \(u_k=(w^1_k,\ldots,w^d_k)^{\mathsf T}\). The projection conclusion makes \(u_k,u_\ell\) linearly independent for \(k\ne\ell\), so every \(u_k\ne0\). Choose \(\beta\in\mathbb R^d\) outside the finitely many proper hyperplanes \(\{\beta:\beta^{\mathsf T}u_k=0\}\). Then choose \(\alpha\in\mathbb R^d\) outside the finitely many proper hyperplanes with normals
\[
(\beta^{\mathsf T}u_\ell)u_k-(\beta^{\mathsf T}u_k)u_\ell,
\qquad k<\ell.
\tag{22}
\]
Every displayed normal is nonzero, because both scalar coefficients are nonzero and \(u_k,u_\ell\) are not proportional. These choices exist: the product of the finitely many nonzero defining linear forms is a nonzero polynomial, and induction on the number of variables using the one-variable root bound proves that a nonzero real polynomial cannot vanish on all of \(\mathbb R^d\).

Set \(b=\sum_j\beta_jw^j\), \(a=\sum_j\alpha_jw^j\), and \(\rho_k=a_k/b_k\). Then every \(b_k\ne0\), and (22) gives \(\rho_k\ne\rho_\ell\) for \(k\ne\ell\). By (21) and linearity, both
\[
G=Q\operatorname{diag}(a)Q^{\mathsf T},
\qquad
K=Q\operatorname{diag}(b)Q^{\mathsf T}
\]
are diagonal, and \(K\) is invertible. Therefore
\[
GK^{-1}=Q\operatorname{diag}(\rho_1,\ldots,\rho_p)Q^{-1}
\tag{23}
\]
is diagonal with pairwise distinct eigenvalues. Its eigenspaces are coordinate axes, so every column of \(Q\) is supported on a coordinate axis. Invertibility makes \(Q\) monomial: there are a permutation \(\pi\) and nonzero scalars \(q_i\) such that
\[
Q_{i,\pi(i)}=q_i,
\qquad Q_{ij}=0\quad(j\ne\pi(i)).
\tag{24}
\]

It remains to remove the permutation at the one-sided normalization boundary. Since \(A=QD\) and both matrices have unit diagonal,
\[
1=A_{ii}=q_iD_{\pi(i),i}.
\tag{25}
\]
If \(\pi\) had a nontrivial cycle \(i_1,\ldots,i_r\), with \(\pi(i_t)=i_{t+1}\) cyclically, (25) would give
\[
\prod_{t=1}^r|D_{i_{t+1},i_t}|
=\prod_{t=1}^r|q_{i_t}|^{-1}
\leq\operatorname{CP}(I_p-D)
\leq1-\zeta_0<1.
\tag{26}
\]
On the reverse cycle, \(A_{i_t,i_{t+1}}=q_{i_t}D_{i_{t+1},i_{t+1}}=q_{i_t}\), so
\[
\prod_{t=1}^r|q_{i_t}|
\leq\operatorname{CP}(I_p-A)
\leq1.
\tag{27}
\]
Multiplying (26) and (27) gives the contradiction \(1<1\). Thus \(\pi\) is the identity; (25) then gives \(q_i=1\) for all \(i\). Hence \(Q=I_p\), or \(AU=I_p\), contradicting (18). Therefore
\[
\mu_0>0.
\tag{28}
\]
Set, non-effectively,
\[
r_0=\min\left\{1,\frac{\mu_0}{2\sqrt p\,L_0^2}\right\}>0.
\tag{29}
\]
The comparison family and its residual use only \(p,m,c,\gamma_0,\zeta_0,\overline\kappa,\overline M\), so \(r_0\) has exactly the asserted dependence. The invoked compactness and minimum-attainment results give no procedure for computing \(\mu_0\) or \(r_0\).

\textbf{Entry into the local neighborhood.}
Fix an outcome in \(G_{n,\alpha,H}\), write \(r=r_{n,\alpha,H}\), and assume \(r\leq r_0\). By \(\text{thm:confidence-union-coverage}\), \(D\in\mathcal U^{(n)}_\alpha\), so the union is nonempty. Let \(A\in\mathcal U^{(n)}_\alpha\). By \(\text{def:confidence-union}\), some simultaneous selection \(\Gamma_e\in\mathcal C^{(n)}_{e,\alpha}\) satisfies \(A\in\mathcal R_c(\Gamma)\). By \(\text{def:compatible-set}\), after restricting a fitted set if necessary, there are \(K\subseteq E\), \(|K|=h\), \(\Psi\in\mathbb S_+^p\), and \(t_e\in[0,\infty)^p\) such that
\[
A\Gamma_eA^{\mathsf T}=\Psi+\operatorname{diag}(t_e)
\quad(e\in K).
\tag{30}
\]
The declarative premise gives the candidate condition-number bound. Membership \(A\in\mathcal D\), through \(\text{def:admissible-matrices}\), gives unit diagonal and \(\operatorname{CP}(I_p-A)<1\), hence the weak closure condition in (8), but no uniform candidate slack is used.

Inclusion--exclusion gives
\[
|H\cap K|\geq|H|+|K|-|E|=2h-m=q.
\tag{31}
\]
Choose \(S\subseteq H\cap K\) with \(|S|=q\), and set \(U=D^{-1}\), \(V=A^{-1}\). For every \(e\in S\), \(\text{def:honest-region-radius}\) gives
\[
\|\Gamma_e-\Sigma_e\|_{\mathrm{op}}
\leq r\leq r_0\leq1.
\tag{32}
\]
The true assumptions and four true-side bounds, (30), the candidate condition-number bound and pointwise membership in \(\mathcal D\), and (32) place this tuple in \(\mathfrak T\). Using (3), (29), (30), and (32), for every \(e\in S\),
\[
\begin{aligned}
\left\|A\Sigma_eA^{\mathsf T}-\Psi-\operatorname{diag}(t_e)\right\|_{\mathrm F}
&=\left\|A(\Sigma_e-\Gamma_e)A^{\mathsf T}\right\|_{\mathrm F}\\
&\leq\sqrt p\,\|A\|_{\mathrm{op}}^2
\|\Sigma_e-\Gamma_e\|_{\mathrm{op}}\\
&\leq\sqrt p\,L_0^2r
\leq\mu_0/2.
\end{aligned}
\tag{33}
\]
Thus the tuple is not in \(\mathfrak T_{\mathrm{far}}\), whether that family is empty or not. With
\[
Q=AD^{-1}=I_p+R,
\qquad u=\|R\|_{\mathrm F},
\]
we obtain
\[
u<\rho_0\leq1,
\qquad K_0\overline M u<\frac12.
\tag{34}
\]

\textbf{Local affine-minor inversion.}
Fix \(e_0\in S\) and write \(x_e=s_e-s_{e_0}\). For every \(i<j\), affine separation on the particular set \(S\) supplies distinct \(a,b,d\in S\) for which the affine determinant in coordinates \((i,j)\) has absolute value at least \(\gamma_0\). With \(\widetilde x_e=(x_{ei},x_{ej})\), translation invariance and direct expansion give
\[
\det(\widetilde x_a-\widetilde x_b,
\widetilde x_d-\widetilde x_b)
=\det(\widetilde x_a,\widetilde x_d)
+\det(\widetilde x_d,\widetilde x_b)
+\det(\widetilde x_b,\widetilde x_a).
\tag{35}
\]
The triangle inequality therefore supplies \(e,f\in S\) such that
\[
|x_{ei}x_{fj}-x_{ej}x_{fi}|\geq\gamma_0/3.
\tag{36}
\]
Let \(E_e=\Gamma_e-\Sigma_e\). Subtract the true equation at \(e_0\) from that at \(e\), do the same in (30), use \(A=QD\), and take the \((i,j)\) off-diagonal entry. The invariant terms cancel, and expansion of \(Q=I_p+R\) yields the exact identity
\[
R_{ij}x_{ej}+R_{ji}x_{ei}
=-\sum_{k=1}^pR_{ik}R_{jk}x_{ek}
-[A(E_e-E_{e_0})A^{\mathsf T}]_{ij}.
\tag{37}
\]
By \(M_H\leq\overline M\) and nonnegativity of the shifts, \(|x_{ek}|\leq\overline M\). Equation (32) gives \(\|E_e-E_{e_0}\|_{\mathrm{op}}\leq2r\), (3) gives \(\|A\|_{\mathrm{op}}\leq L_0\), and each row of \(R\) has Euclidean norm at most \(u\). Cauchy--Schwarz in (37) therefore gives
\[
|R_{ij}x_{ej}+R_{ji}x_{ei}|
\leq\overline M u^2+2L_0^2r.
\tag{38}
\]
Apply (38) at the two labels in (36). The determinant of the resulting two-by-two coefficient matrix for \((R_{ij},R_{ji})\) has absolute value at least \(\gamma_0/3\), every coefficient has absolute value at most \(\overline M\), and each right-hand side has absolute value at most \(\beta=\overline M u^2+2L_0^2r\). Cramer's rule and the triangle inequality give, for every ordered pair \(i\ne j\),
\[
|R_{ij}|
\leq\frac{6\overline M}{\gamma_0}
\left(\overline M u^2+2L_0^2r\right).
\tag{39}
\]
Because \(A\) and \(D\) have unit diagonal and \(A-D=RD\),
\[
R_{ii}=-\sum_{k\ne i}R_{ik}D_{ki}.
\tag{40}
\]
The squared Euclidean norm of every column of \(D\) is at most \(L_0^2\) by (3). Squaring (39), summing over the \(p(p-1)\) ordered off-diagonal entries, applying Cauchy--Schwarz to (40), and taking square roots yields
\[
u\leq K_0\left(\overline M u^2+2L_0^2r\right).
\tag{41}
\]
If \(u=0\), the next inequality is immediate. If \(u>0\), (34) makes \(K_0\overline M u^2<u/2\); subtracting this term from (41) gives
\[
u\leq4K_0L_0^2r.
\tag{42}
\]
Using (3), (42), and \(C_0=16K_0L_0^3\),
\[
\|A-D\|_{\mathrm{op}}
=\|RD\|_{\mathrm{op}}
\leq uL_0
\leq4K_0L_0^3r
\leq C_0r.
\tag{43}
\]
Since \(A\in\mathcal U^{(n)}_\alpha\) was arbitrary, (43) and \(\text{def:operator-outer-radius}\) give the outer-radius bound. For \(A_1,A_2\in\mathcal U^{(n)}_\alpha\), the triangle inequality and (43) give
\[
\|A_1-A_2\|_{\mathrm{op}}
\leq\|A_1-D\|_{\mathrm{op}}+\|A_2-D\|_{\mathrm{op}}
\leq2C_0r.
\]
Taking the supremum and applying \(\text{def:operator-diameter}\) proves the diameter bound. Along any covered sequence in the statement, \(r_{n,\alpha,H}\to0\) implies eventually \(r_{n,\alpha,H}\leq r_0\); the two bounds then imply both limits. No candidate affine separation, candidate scale bound, uniform candidate normalization slack, or covariance-interiority assumption has been used.